\chapter{General Renormalization}
\label{ch:renormalization}

Coarse arrows alone form a scale diagram. They become a renormalization
dynamics only after comparison and rescaling data identify what is being
compared across levels. In the present theory the state includes one
principal gauge bundle, two associated statistical bundles, their possibly
distinct connections, and both cross-bundle morphisms. A scalar normalization
of one operator cannot stand for that full state.

\section{The scale category and comparison data}
\label{sec:rg-scale-category}

\definitionheading{Scale diagram}{def:rg-scale-diagram}
Let \(\mathsf S\) be a category whose objects are resolutions \(\ell\).
Choose a target category \(\mathscr K\): sets and maps for a deterministic
state, standard Borel spaces and Markov kernels for a stochastic state, an
appropriate linear category for an operator state, or an explicitly declared
product of such categories for a mixed state. A state functor
\(F:\mathsf S\to\mathscr K\) assigns an object
\(\mathfrak X_\ell=F(\ell)\) and morphisms
\begin{equation}
C_{\ell k}\in
\operatorname{Hom}_{\mathscr K}(\mathfrak X_\ell,\mathfrak X_k),
\qquad
C_{kr}C_{\ell k}=C_{\ell r}.
\label{eq:rg-scale-diagram}
\end{equation}
The arrow denotes the appropriate composition in \(\mathscr K\), including
kernel composition in the stochastic case and componentwise composition in a
product category. One symbol is never asked to be simultaneously a map, a
kernel, and an operator. \status{DEFINITION}

Comparison requires more data. Either choose levelwise re-embeddings
\(\jmath_\ell:\mathfrak X_{\ell+1}\to\mathfrak X_\ell\), or identify each
level with a reference object by declared isomorphisms
\(I_\ell:\mathfrak X_\ell\xrightarrow{\sim}\mathfrak X_\star\) in
\(\mathscr K\).
Field, base, unit, and reference-measure rescalings are part of these maps.
If an externally parameterized dynamical model is separately supplied, its
parameter rescaling also belongs here; the scale diagram itself introduces no
time variable, and RG depth is not the intrinsic inference duration of
\Cref{ch:relational-inference}. Only the isomorphism option supplies the inverse used in the
reference-space transformations
\begin{equation}
\widehat{\mathcal R}_\ell
=I_{\ell+1}C_{\ell,\ell+1}I_\ell^{-1}:
\mathfrak X_\star\longrightarrow\mathfrak X_\star.
\label{eq:rg-reference-cocycle}
\end{equation}
\status{DEFINITION} A positive scalar \(\zeta_\ell\) can be one component of
this comparison; it cannot by itself identify changing fields, measures, or
bundle objects.

If probability laws are present, every scale declares its sample space and
reference measure. A law-level arrow is a normalized Markov kernel as in
\Cref{ch:coarsegraining}. An energy-level arrow additionally owes the
coarse reference measure, its Radon--Nikodym or Jacobian rule, and a finite
normalizer. Algebraic closure does not supply those data. \status{HYPOTHESIS}

\section{Exact measure-pair arrows}
\label{sec:rg-measure-pair-arrows}

\definitionheading{Scale-indexed measure pair}{def:rg-measure-pair}
For each scale \(\ell\), let \((\mathsf X_\ell,\mathscr X_\ell)\) be a
standard-Borel space, let \(\rho_\ell\) be a probability measure, and let
\(m_\ell\ll\rho_\ell\) be a finite positive measure.  Write
\begin{equation}
m_\ell(dx)=\exp[-H_\ell(x)]\rho_\ell(dx),
\qquad
0<M_\ell:=m_\ell(\mathsf X_\ell)<\infty,
\qquad
\pi_\ell=\frac{m_\ell}{M_\ell}.
\label{eq:rg-measure-pair-normalization}
\end{equation}
The extended-value convention is \(\exp(-\infty)=0\); thus the action is
defined only \(\rho_\ell\)-almost everywhere and no finite action is
silently assumed.  A scale arrow is a normalized, parameter-independent
Markov kernel
\begin{equation}
K_\ell:\mathsf X_\ell\rightsquigarrow\mathsf X_{\ell+1},
\qquad
\rho_{\ell+1}=\rho_\ell K_\ell,
\qquad
m_{\ell+1}=m_\ell K_\ell.
\label{eq:rg-measure-pair-arrow}
\end{equation}
The downstream action is the Radon--Nikodym representative
\begin{equation}
H_{\ell+1}=-\log\frac{dm_{\ell+1}}{d\rho_{\ell+1}}
\quad \rho_{\ell+1}\text{-almost everywhere}.
\label{eq:rg-measure-pair-action}
\end{equation}
This is a definition of an exact law-level RG arrow.  The requirement that
the spaces be standard Borel is used later to choose regular conditional
versions; it does not identify distinct scale spaces or produce an
autonomous operator. \status{DEFINITION}

\propositionheading{Exact composition and normalized law}{prop:rg-measure-pair-composition}
For two composable normalized kernels \(K_\ell\) and \(K_{\ell+1}\), the
measure-pair arrows compose exactly:
\begin{equation}
(\rho_\ell,m_\ell)(K_\ell K_{\ell+1})
=\bigl((\rho_\ell K_\ell)K_{\ell+1},
(m_\ell K_\ell)K_{\ell+1}\bigr).
\label{eq:rg-measure-pair-composition}
\end{equation}
Moreover \(m_{\ell+1}\ll\rho_{\ell+1}\),
\(m_{\ell+1}(\mathsf X_{\ell+1})=M_\ell\), and
\begin{equation}
\pi_{\ell+1}=\pi_\ell K_\ell.
\label{eq:rg-normalized-law-pushforward}
\end{equation}
Thus the normalized law and the unnormalized evidence mass are distinct
components of the same exact construction. \status{ESTABLISHED}

\paragraph{Proof.}
Kernel composition and Tonelli's theorem give
\eqref{eq:rg-measure-pair-composition}.  If
\((\rho_\ell K_\ell)(A)=0\), then
\(K_\ell(x,A)=0\) for \(\rho_\ell\)-almost every \(x\), hence also for
\(m_\ell\)-almost every \(x\), which proves absolute continuity at the
next scale.  Normalization of \(K_\ell\) preserves the mass of
\(m_\ell\); division by that same mass gives
\eqref{eq:rg-normalized-law-pushforward}.  \(\square\)

The subsequent finite-network realization takes each
\(\mathsf X_\ell\) to be a finite product of agent coordinate spaces and
supplies the action calculus, contraction theorems, and interaction content
that this general arrow does not itself supply.  The exact nonlinear
interaction-coordinate map, its tilted derivative, and its retained-space
residual are defined only after the additional product-equivalence premises
of \Cref{sec:rg-finite-network-interactions}; they are not linear
conditional-expectation maps away from the unperturbed action.  In particular,
a scale sequence is a cocycle until the comparison data of
\Cref{sec:rg-scale-category} identify its varying spaces. \status{ESTABLISHED}

\section{The common-principal geometric state}
\label{sec:rg-common-principal-state}

\definitionheading{Geometric RG state}{def:rg-geometric-state}
At level \(\ell\), declare a smooth base \(\mathcal C_\ell\), a Lie group
\(G_\ell\), and one principal right \(G_\ell\)-bundle
\(\pi_\ell:P_\ell\to\mathcal C_\ell\). Let
\(\mathcal B_{\ell,b}\) and \(\mathcal B_{\ell,m}\) be statistical fibers with
possibly different left actions
\(\widehat\rho_{\ell,b}\) and \(\widehat\rho_{\ell,m}\) of the same group. The
belief and model bundles are
\begin{equation}
\mathcal E_{\ell,b}
=P_\ell\times_{\widehat\rho_{\ell,b}}\mathcal B_{\ell,b},
\qquad
\mathcal E_{\ell,m}
=P_\ell\times_{\widehat\rho_{\ell,m}}\mathcal B_{\ell,m}.
\label{eq:rg-associated-bundles}
\end{equation}
Choose two principal connections
\(\omega_{\ell,b},\omega_{\ell,m}\in
\Omega^1(P_\ell,\mathfrak g_\ell)\). They induce horizontal distributions
\(H^b_\ell,H^m_\ell\) and parallel transports
\(\Omega_{\ell,\gamma}\) and
\(\widetilde\Omega_{\ell,\gamma}\) on the two associated bundles. Thus the
general geometric state contains
\begin{equation}
\mathfrak X_\ell=
\left(
\mathcal C_\ell;\
G_\ell,P_\ell;\
\mathcal E_{\ell,b},H^b_\ell;\
\mathcal E_{\ell,m},H^m_\ell;\
\Phi_\ell,\widetilde\Phi_\ell;\
\ldots
\right).
\label{eq:rg-full-state}
\end{equation}
Here
\[
\Phi_\ell:\mathcal E_{\ell,b}\to\mathcal E_{\ell,m},
\qquad
\widetilde\Phi_\ell:\mathcal E_{\ell,m}\to\mathcal E_{\ell,b}
\]
are separately declared associated-bundle morphisms over the identity. They
need not be invertible or mutual inverses. A fixed fiber map
\(f^{bm}_\ell:\mathcal B_{\ell,b}\to\mathcal B_{\ell,m}\) induces such a morphism
only when it intertwines the actions,
\begin{equation}
f^{bm}_\ell\circ\widehat\rho_{\ell,b}(g)
=\widehat\rho_{\ell,m}(g)\circ f^{bm}_\ell
\quad\text{for every }g\in G_\ell,
\label{eq:rg-cross-intertwiner}
\end{equation}
and analogously in the reverse direction. Hence a common principal bundle
does not manufacture either cross map. The ellipsis in
\eqref{eq:rg-full-state} denotes law, kernel, or operator components declared
by a realization. \status{DEFINITION}

When levelwise agent or meta-agent sections and regular vertical Fisher
tensors are also declared, the geometric state may include their
connection-relative perceived semimetrics.  These components are not supplied
by the associated bundles alone.  If a scale arrow also supplies a bundle
morphism, related fine and coarse sections, and a vanishing horizontal
defect, their exact cross-scale naturality is
\Cref{cor:pb-meta-perceived-geometry}.  Its positive-semidefinite
information-loss defect further requires that the fiber map be a normalized
parameter-independent Markov channel.  The discrete or continuous beta is
defined only after the reference-base identifications in
\eqref{eq:pb-metric-beta-reference}.  Thus a meta-agent has its own pullback
geometry, but no tensor on two different bases is subtracted before transport
to the common reference space. \status{ESTABLISHED}

Local belief and model frames may nevertheless be chosen separately. On a
common trivializing neighborhood, two sections
\(u^b_{\ell,i},u^m_{\ell,i}:U_i\to P_\ell\) determine a unique smooth relative
principal-frame field
\begin{equation}
u^m_{\ell,i}=u^b_{\ell,i}\cdot h_{\ell,i},
\qquad h_{\ell,i}:U_i\to G_\ell.
\label{eq:rg-relative-frame}
\end{equation}
Existence and uniqueness follow because each principal fiber is a free and
transitive \(G_\ell\)-torsor; the statement is independent of the
representations. If the sections are
re-chosen as \(u^b_{\ell,i}a^b_{\ell,i}\) and
\(u^m_{\ell,i}a^m_{\ell,i}\), then
\(h_{\ell,i}\mapsto
(a^b_{\ell,i})^{-1}h_{\ell,i}a^m_{\ell,i}\). When the representations
differ, \(\widehat\rho_{\ell,b}(h_{\ell,i})\) and
\(\widehat\rho_{\ell,m}(h_{\ell,i})\) act on different fibers; there is no
matrix quotient of the associated frames and no cross-fiber identification
without extra intertwining data such as \eqref{eq:rg-cross-intertwiner}.
If only associated-frame data are given and either action is nonfaithful, the
principal element \(h_{\ell,i}\) need not be reconstructible without chosen
lifts to \(P_\ell\).
\status{ESTABLISHED}

Scale change must also be typed. Let
\(c_\ell:\mathcal C_\ell\to\mathcal C_{\ell+1}\) be smooth, let
\(\kappa_\ell:G_\ell\to G_{\ell+1}\) be a Lie-group homomorphism, and let
\(\mathcal P_\ell:P_\ell\to P_{\ell+1}\) cover \(c_\ell\) and obey
\begin{equation}
\mathcal P_\ell(p\cdot g)
=\mathcal P_\ell(p)\cdot\kappa_\ell(g).
\label{eq:rg-principal-scale-map}
\end{equation}
For \(s\in\{b,m\}\), a fiber map
\(q_{\ell,s}:\mathcal B_{\ell,s}\to\mathcal B_{\ell+1,s}\) descends to a
bundle map only if
\begin{equation}
q_{\ell,s}\circ\widehat\rho_{\ell,s}(g)
=\widehat\rho_{\ell+1,s}(\kappa_\ell(g))\circ q_{\ell,s}.
\label{eq:rg-scale-intertwiner}
\end{equation}
It then induces
\begin{equation}
C_{\ell,s}:\mathcal E_{\ell,s}\to\mathcal E_{\ell+1,s},
\qquad
C_{\ell,s}[p,z]=[\mathcal P_\ell(p),q_{\ell,s}(z)],
\label{eq:rg-associated-scale-map}
\end{equation}
covering \(c_\ell\). More general associated-bundle maps may be declared
directly, but they owe the same transition compatibility. For a genuine
multi-step scale functor, the maps \(c\), \(\kappa\), \(\mathcal P\), and
\(q_{\ell,s}\) also obey their respective composition laws; adjacent maps without
these laws form only a sequence of coarse steps. If local comparison functions
\(a^s_\ell\) satisfy
\[
\mathcal P_\ell(u^s_\ell(x))
=u^s_{\ell+1}(c_\ell x)\cdot a^s_\ell(x),
\]
then equivariance of \(\mathcal P_\ell\) forces
\begin{equation}
a^b_\ell\kappa_\ell(h_\ell)
=(h_{\ell+1}\circ c_\ell)a^m_\ell.
\label{eq:rg-relative-frame-naturality}
\end{equation}
Thus the two represented frame channels remain compatible even when their
fiber dimensions or actions differ. \status{ESTABLISHED}

For general nonlinear associated fibers, the cross-scale morphism defects are
the ordered pairs of maps
\begin{equation}
\mathbf D^\Phi_\ell
=\left(C_{\ell,m}\circ\Phi_\ell,
\Phi_{\ell+1}\circ C_{\ell,b}\right),\qquad
\mathbf D^{\widetilde\Phi}_\ell
=\left(C_{\ell,b}\circ\widetilde\Phi_\ell,
\widetilde\Phi_{\ell+1}\circ C_{\ell,m}\right).
\label{eq:rg-cross-morphism-defects}
\end{equation}
The defect vanishes when the two components of the corresponding pair are
equal. For a base curve \(\gamma:x\to y\), the cross-scale connection
defects are likewise the pairs
\begin{align}
\mathbf D^b_{\ell,\gamma}
&=\left(C_{\ell,b,y}\circ\Omega_{\ell,\gamma},
\Omega_{\ell+1,c_\ell\gamma}\circ C_{\ell,b,x}\right),\nonumber\\
\mathbf D^m_{\ell,\gamma}
&=\left(C_{\ell,m,y}\circ\widetilde\Omega_{\ell,\gamma},
\widetilde\Omega_{\ell+1,c_\ell\gamma}\circ C_{\ell,m,x}\right).
\label{eq:rg-cross-connection-defects}
\end{align}
Independently, within-scale nonparallelness is measured by
\begin{align}
\mathbf A^\Phi_{\ell,\gamma}
&=\left(\widetilde\Omega_{\ell,\gamma}\circ\Phi_{\ell,x},
\Phi_{\ell,y}\circ\Omega_{\ell,\gamma}\right),\nonumber\\
\mathbf A^{\widetilde\Phi}_{\ell,\gamma}
&=\left(\Omega_{\ell,\gamma}\circ\widetilde\Phi_{\ell,x},
\widetilde\Phi_{\ell,y}\circ\widetilde\Omega_{\ell,\gamma}\right).
\label{eq:rg-within-connection-defects}
\end{align}
\status{DEFINITION} Subtraction of the two components is reserved for a
declared vector or affine realization. A sufficient principal-level condition
for \(\mathbf D^s_{\ell,\gamma}\) to vanish is
\begin{equation}
\mathcal P_\ell^*\omega_{\ell+1,s}
=\mathrm d\kappa_\ell\circ\omega_{\ell,s},
\qquad s\in\{b,m\}.
\label{eq:rg-principal-connection-naturality}
\end{equation}
If a representation is not faithful, vanishing of its represented transport
defect need not imply this principal-level identity.

\propositionheading{Covariance of the defect diagrams}{prop:rg-defect-covariance}
Under admissible local-frame changes in the belief and model associated
bundles, both components of every pair in
\eqref{eq:rg-cross-morphism-defects}--\eqref{eq:rg-within-connection-defects}
acquire the same endpoint coordinate changes. Consequently the equality
locus of each pair is gauge invariant. In a fiberwise-linear vector-bundle
realization, subtracting the two components gives the familiar Hom-bundle
defect with the same left--right law. \status{ESTABLISHED}

\paragraph{Proof.}
Let \(\mathscr G_{\ell,b,x}\) and \(\mathscr G_{\ell,m,x}\) denote the fiber
coordinate changes induced by the selected belief and model frame
re-trivializations at \(x\), represented through
\(\widehat\rho_{\ell,b}\) and \(\widehat\rho_{\ell,m}\). For example,
\[
\Phi_{\ell,x}\mapsto
\mathscr G_{\ell,m,x}\circ\Phi_{\ell,x}\circ
(\mathscr G_{\ell,b,x})^{-1},
\quad
\Omega_{\ell,\gamma}\mapsto
\mathscr G_{\ell,b,y}\circ\Omega_{\ell,\gamma}\circ
(\mathscr G_{\ell,b,x})^{-1},
\]
and a cross-scale map obeys
\[
C_{\ell,b,x}\mapsto
\mathscr G_{\ell+1,b,c_\ell(x)}\circ C_{\ell,b,x}\circ
(\mathscr G_{\ell,b,x})^{-1},
\]
with analogous model and endpoint laws. Substitution shows that both
components of \(\mathbf A^\Phi_{\ell,\gamma}\) acquire the same model
coordinate change at \(y\) and inverse belief coordinate change at \(x\).
The equality of the two components is therefore independent of the chosen
frames. The other pairs are identical calculations. In a linear realization,
subtraction factors out the common endpoint actions. \(\square\)

The defect data may be retained as running fields; the theory does not
require them to vanish. A numerical size requires a declared metric or
divergence on each target fiber. The compositions
\(\widetilde\Phi_\ell\Phi_\ell\) and
\(\Phi_\ell\widetilde\Phi_\ell\) are independent running endomorphisms and
are not identities unless imposed. Smooth parallel transport remains
distinct from separately declared graph links
\(\Theta^b_{ij},\Theta^m_{ij}\). \status{ESTABLISHED}

The default common-principal construction does not impose a common local
frame or a common connection. The stronger common-frame hypothesis is
\(u^b_{\ell,i}=u^m_{\ell,i}\), equivalently
\(h_{\ell,i}=e\), on the chosen patches. Independently, the equal-connection
hypothesis is \(\omega_{\ell,b}=\omega_{\ell,m}\). Even then, different
representations generally produce different associated transports on
different fibers: they represent the same principal holonomy but are not
literally the same map. If a further associated-bundle isomorphism
\(\Xi_\ell:\mathcal E_{\ell,b}\xrightarrow{\sim}\mathcal E_{\ell,m}\) is
declared, its existence is additional global data; equal fiber type is
necessary but not sufficient in general. Its parallel and RG-preservation
conditions are
\[
\widetilde\Omega_{\ell,\gamma}\Xi_{\ell,x}
=\Xi_{\ell,y}\Omega_{\ell,\gamma},\qquad
C_{\ell,m}\Xi_\ell=\Xi_{\ell+1}C_{\ell,b}.
\]
The optional \(\Xi_\ell\) is unavailable when the associated bundles are
nonisomorphic, as can occur when ranks or fiber types differ. No such
isomorphism is needed to define \(\Phi_\ell\) and
\(\widetilde\Phi_\ell\). \status{HYPOTHESIS}

\paragraph{Optional product-gauge extension.}
If an application truly requires independent gauge groups, independent
principal-bundle topology, or independently acting gauge automorphisms, one
may instead introduce
\[
P_{\ell,b}\to\mathcal C_\ell,\qquad
P_{\ell,m}\to\mathcal C_\ell,\qquad
Q_\ell=P_{\ell,b}\times_{\mathcal C_\ell}P_{\ell,m},
\]
where \(Q_\ell\) is principal for
\(G_{\ell,b}\times G_{\ell,m}\). Each associated bundle then uses its own
factor. This extension carries extra topological and gauge data and is not
part of the minimal theory. In a category admitting the required mapping
object, its belief-to-model maps are sections of
\[
Q_\ell\times_{G_{\ell,b}\times G_{\ell,m}}
\operatorname{Map}(\mathcal B_{\ell,b},\mathcal B_{\ell,m}),
\qquad
(g_b,g_m)\cdot f
=\widehat\rho_{\ell,m}(g_m)\circ f\circ
\widehat\rho_{\ell,b}(g_b)^{-1}.
\]
A constant fiber map descends only if it is fixed by this action, a generally
restrictive condition for independently varying \((g_b,g_m)\). Cross maps
therefore still have to be supplied, and a diagonal reduction to one group
requires additional group and bundle-identification data. \status{DEFINITION}

\section{Autonomous maps and cocycles}
\label{sec:rg-dynamical-types}

\definitionheading{Dynamical type}{def:rg-dynamical-type}
The sequence \(\widehat{\mathcal R}_\ell\) in
\eqref{eq:rg-reference-cocycle} is:
autonomous when it is one map \(\widehat{\mathcal R}\);
periodic when
\(\widehat{\mathcal R}_{\ell+p}=\widehat{\mathcal R}_\ell\);
stationary random when it is a measurable stationary cocycle; and general
otherwise. \status{DEFINITION}

The correct invariant object depends on this type. An autonomous fixed object
satisfies \(\widehat{\mathcal R}x=x\). A periodic scheme has a fixed object
of its monodromy or an attracting \(p\)-cycle. A stationary random scheme may
have a random invariant section. A general cocycle may have pullback or
forward attracting sections. The term \emph{fixed ray} additionally requires
a positive cone and projectivization. \status{DEFINITION}

\propositionheading{Uniform contraction need not give one fixed ray}{prop:rg-contraction-no-fixed-ray}
Uniform Hilbert-projective contraction of a nonautonomous positive cocycle can
erase initial conditions without selecting one scale-independent ray.
\status{ESTABLISHED}

\paragraph{Proof.}
On \(\R^2_{>0}\), let
\[
A=\begin{pmatrix}2&1\\1&1\end{pmatrix},
\qquad
B=\begin{pmatrix}1&1\\1&2\end{pmatrix}.
\]
Both are strictly positive and have a uniform projective contraction
coefficient below one. Their Perron rays differ. Under the alternating
cocycle, even iterates converge to the Perron ray
\([1,\sqrt2]\) of
\[
BA=\begin{pmatrix}3&2\\4&3\end{pmatrix},
\]
while odd iterates converge to the distinct ray
\(A[1,\sqrt2]\). The attractor is period two, not one fixed ray. \(\square\)

For maps satisfying
\(\tau(\widehat{\mathcal R}_\ell)\leq q<1\), Birkhoff contraction gives
\begin{equation}
d_H(
\widehat{\mathcal R}_{\ell-1}\cdots\widehat{\mathcal R}_0x,
\widehat{\mathcal R}_{\ell-1}\cdots\widehat{\mathcal R}_0y)
\leq q^\ell d_H(x,y),
\label{eq:rg-cocycle-contraction}
\end{equation}
where defined \citep{Birkhoff1957,nussbaum1988hilbert}. This proves forgetting
of initial rays, not time independence of the attracting section.

\section{Scheme dependence and invariant candidates}
\label{sec:rg-scheme-dependence}

Let \(g(t)\) be running coordinates. Under a smooth reparameterization
\(g'=f(g)\) and scale change \(t'=ct\),
\begin{equation}
\beta'(g')=\frac1cDf(g)\beta(g).
\label{eq:rg-beta-change}
\end{equation}
\status{ESTABLISHED} Hence raw couplings, beta-function components, and a
chosen blocking statistic are scheme dependent. They remain valid diagnostics
when their scheme is stated, but are not universal observables.

For a differentiable homogeneous positive endomorphism at a fixed ray, let
\(\mu_0\) be the radial eigenvalue and \(\mu_a\) a transverse eigenvalue.
The projective ratios and block-normalized dimensions are
\begin{equation}
\rho_a=\frac{\mu_a}{\mu_0},\qquad
y_a=\frac{\log|\rho_a|}{\log b},
\label{eq:rg-projective-dimensions}
\end{equation}
for declared block factor \(b\). \status{ESTABLISHED} They are invariant
under smooth coordinate conjugacy at the ray once gauge, normalization, and
other redundant directions have been quotiented and the same scale
normalization is used.

Candidate cross-scheme quantities include conjugacy classes, invariant
manifolds, projective eigenvalue ratios, dimensionless amplitude ratios, and
matched spectral or heat exponents. Their existence, convergence, and
completeness must be proved in the chosen realization. Universality is a
statement about basins and invariant observables, not algebraic closure
alone. \status{HYPOTHESIS}

\section{Two-index thermodynamic and RG limits}
\label{sec:rg-two-index-limits}

A finite-volume RG family has two indices:
\begin{equation}
X_{n,\ell},\qquad 0\leq\ell\leq L(n),
\label{eq:rg-two-index-family}
\end{equation}
where \(n\) controls system size and \(\ell\) RG depth. Fixed-depth,
maximal-depth, and diagonal limits are different:
\begin{equation}
\lim_{\ell\to\infty}\lim_{n\to\infty}X_{n,\ell},
\qquad
\lim_{n\to\infty}X_{n,L(n)},
\qquad
\lim_{n\to\infty}X_{n,\ell(n)}.
\label{eq:rg-two-index-limits}
\end{equation}
\status{DEFINITION}

\propositionheading{The two limits need not commute}{prop:rg-noncommuting-limits}
For path graphs with hard blocks of size \(b\), one coarse step maps
\(P_{b^n}\) to \(P_{b^{n-1}}\), up to endpoint convention. At each fixed
\(\ell\), \(n\to\infty\) yields the infinite path, whose lower spectral edge
has counting behavior \(N(\lambda)-N(0)\asymp\lambda^{1/2}\). Along
\(\ell=n\), every graph collapses to one vertex and the spectral measure is
\(\delta_0\). \status{ESTABLISHED}

\paragraph{Proof.}
Hard aggregation of consecutive path vertices preserves the path incidence
pattern. At fixed \(\ell\), the number of remaining vertices tends to
infinity; at \(\ell=n\), it is one. The finite path eigenvalues have the
standard small-wave-number quadratic behavior, yielding the square-root
counting edge. The one-vertex operator is zero. \(\square\)

Therefore a thermodynamic RG claim must specify a triangular family,
projective consistency at fixed depth, uniform control of the coarse error,
and a diagonal scaling \(\ell(n)\). Operator, law, connection, and
cross-morphism convergence are separate obligations. Proving an exchange of
limits or a nontrivial diagonal limit remains open. \status{OPEN}

If an infinite-volume integrated density satisfies
\[
N(\lambda)-N(0)\sim c\lambda^\alpha
\]
and a Tauberian hypothesis supplies the corresponding heat asymptotic, then
the nonzero heat trace and unnormalized spectral entropy obey
\begin{equation}
Z(t)\sim c\Gamma(\alpha+1)t^{-\alpha},\qquad
-\frac{dS}{d\log t}
=t^2\operatorname{Var}_t(\lambda)\longrightarrow\alpha.
\label{eq:rg-heat-susceptibility}
\end{equation}
\status{ESTABLISHED} Dividing entropy by \(\log|V_n|\) divides the
susceptibility by the same factor and generally destroys a nonzero
thermodynamic plateau. Zero-mode atoms and normalization must therefore be
declared before interpreting a plateau as spectral dimension.

\section{Projective laws and refinement}
\label{sec:rg-projective-refinement}

Compatible finite-dimensional laws obey the projective theorem of
\Cref{sec:cg-projective-laws}. An energy recursion alone does not provide
those laws. A probabilistic RG needs normalized kernels
\(K_{\ell k}\) with
\[
\mu_k=\mu_\ell K_{\ell k},\qquad
K_{\ell k}K_{kr}=K_{\ell r},
\]
and corresponding rules for the two bundle sectors and their reference
measures. Smooth-section support, continuum ELBO convergence, and connection
convergence remain additional open obligations. \status{OPEN}

If the coarse map \(Q\) is noninjective, no deterministic two-sided inverse
exists. A refinement kernel \(R\) can at most be a stochastic right section,
\begin{equation}
Q_\#R(c,\cdot)=\delta_c,
\label{eq:rg-stochastic-right-section}
\end{equation}
which redraws a compatible microstate. It cannot recover every discarded
realized microstate. \status{ESTABLISHED} Infinite divisibility may construct
such refinements in distribution, but does not turn a many-to-one
coarse-graining into a group inverse.

\section{Typed comparison with Bayesian renormalization}
\label{sec:rg-bayesian-comparison}

Berman, Klinger, and Stapleton build on the dynamical Bayesian-inference
flow of \citet{BermanHeckmanKlinger2022}, using a centered-KL posterior flow
on model-parameter space. In the regular late-data regime they approximate the
posterior by a Gaussian with covariance \(T^{-1}I^{-1}\). Set
\(\tau=1/T\), and consider an inverse problem with a differentiable forward
map \(G\). The delta method gives the leading local predictive approximation
\[
G(\theta)\ \dot\sim\
\mathcal N\left(
G(\mu_\tau),\quad
\tau J_G(\mu_\tau) I(\mu_\tau)^{-1}J_G(\mu_\tau)^\top
\right).
\]
It is exact when \(G\) is affine, but a nonlinear Gaussian pushforward is not
in general Gaussian. Accordingly, the pushed inverse Fisher cometric
\begin{equation}
D_\tau
=J_G(\mu_\tau)I(\mu_\tau)^{-1}J_G(\mu_\tau)^\top
\label{eq:rg-bks-diffusion}
\end{equation}
is the local diffusion tensor used in their ERG-form Fokker--Planck equation;
the corresponding small-\(\tau\) transition covariance is
\(\tau D_\tau\) \citep{BermanKlingerStapleton2023}. Their continuation of
this tensor beyond the late-data, small-\(\tau\) regime is a renormalization
construction, not an exact nonlinear pushforward theorem or a Bayesian
asymptotic theorem. \status{ESTABLISHED}

This is an exact, carefully scoped literature statement. It does not identify
an arbitrary recognition Fisher tensor with the BKS generative-model Fisher
metric, transfer their stiff/sloppy eigenspaces to another operator sector, or
prove that a graph aggregation runs in the opposite direction.
\status{NOT-CLAIMED} In the small-\(\tau\) approximation, covariance broadens
like \(\tau I^{-1}\) if \(I\) is held fixed, whereas a later Gaussian
sufficient-statistic construction may add precision. That is a typed local
contrast. A theorem of opposite flow directions would require a common state
space, an explicit scale map, and a transported monotone including
\[
\frac{d}{d\tau}\left[\tau I(\mu_\tau)^{-1}\right]
=I^{-1}-\tau I^{-1}\frac{dI}{d\tau}I^{-1}.
\]
No such bridge is asserted. \status{OPEN}

\Cref{thm:cg-holonomy-kl-marginal} gives a separate
admissibility statement: represented holonomy restricts the marginal-law
modes available inside an already proposed block.  The BKS Fisher scale
instead orders model-parameter directions by statistical distinguishability;
it neither chooses node partitions nor tests holonomy. \status{ESTABLISHED}
Constructing a nondegenerate partition selector compatible with these
separate structures remains open. \status{OPEN}

Mehta and Schwab give an exact variational-RG/deep-network correspondence
under their RBM kernel and trace condition \citep{MehtaSchwab2014}. It is a
historical exact precedent inside that construction, not a theorem that
deterministic energy precomposition is a neural layer. Likewise, scalar
Laplacian RG provides a heat-spectrum scale diagnostic and a heuristic
real-space blocker; lifting its entrywise blocker covariantly to the two
represented frame channels requires a new invariant statistic. These are
constraints on applications, not definitions of the present general RG.
\status{ESTABLISHED}

\paragraph{Physicist summary.}
Blocking removes detail; renormalization adds the maps that say how the
blocked theory is compared with the original units. Here that comparison must
carry one principal bundle, its two differently represented statistical
bundles, two connection choices that may coincide, and two separately
supplied cross morphisms. Separate belief and model frames are local views of the same
principal geometry; their relative element lies in \(G_\ell\), even when its
two representations act on unlike spaces. Failures to commute with blocking
are covariant defect fields. Once all scales are identified, the dynamics may
be autonomous, periodic, random, or a general cocycle, and each case has a
different invariant object. System size and RG depth then form a two-index
limit. Fixed points, universality, and continuum laws are meaningful only
after these types and limits have been declared.
